% =========================================================
%  config/environments.tex
%  Kotak dan lingkungan yang dipakai berulang pada tiap sesi.
% =========================================================

\tcbset{
  sesibox/.style={
    enhanced, breakable, arc=1.5pt, boxrule=0.6pt,
    left=8pt, right=8pt, top=6pt, bottom=6pt,
    fonttitle=\sffamily\bfseries\small, coltitle=white,
    toptitle=2pt, bottomtitle=2pt,
    before skip=8pt, after skip=8pt,
    before upper={\setlength{\parskip}{0.35em}}
  }
}

% --- Tujuan sesi ------------------------------------------
\newtcolorbox{tujuanbox}[1][]{sesibox, colframe=ubNavy, colback=ubSoft,
  title={Yang Dapat Anda Tunjukkan}, #1}

% --- Prasyarat --------------------------------------------
\newtcolorbox{prasyaratbox}[1][]{sesibox, colframe=ubLine, colback=ubGray,
  coltitle=ubNavy, title={Sebelum Memulai}, #1}

% --- Uji hasil di tengah tutorial -------------------------
\newtcolorbox{ujibox}[1][]{sesibox, colframe=ubGreen, colback=ubGreenBg,
  title={Periksa Hasilnya}, #1}

% --- Kesalahan umum ---------------------------------------
\newtcolorbox{galatbox}[1][]{sesibox, colframe=ubRed, colback=ubRedBg,
  title={Gejala, Penyebab, dan Letak Pemeriksaan}, #1}

% --- Lembar kerja -----------------------------------------
\newtcolorbox{lembarbox}[1][]{sesibox, colframe=ubAccent, colback=white,
  title={Isi dan Kumpulkan}, #1}

% --- Catatan ----------------------------------------------
\newtcolorbox{catatanbox}[1][]{sesibox, colframe=ubAccent, colback=ubSoft,
  title={Catatan}, #1}

% --- Checklist sesi ---------------------------------------
\newtcolorbox{checklistbox}[1][]{sesibox, colframe=ubGreen, colback=white,
  title={Centang Bila Sudah Dapat}, #1}

% --- Langkah bernomor -------------------------------------
\newenvironment{langkah}{%
  \begin{tcolorbox}[enhanced, breakable, arc=1.5pt, boxrule=0.6pt,
    colframe=ubHead, colback=white, leftrule=2.5pt, colframe=ubAccent,
    left=8pt, right=8pt, top=6pt, bottom=6pt,
    before skip=6pt, after skip=6pt]
  \begin{enumerate}[label=\textbf{\color{ubAccent}\arabic*.},
    leftmargin=1.55em, itemsep=0.45em, topsep=0pt, parsep=0pt]
}{%
  \end{enumerate}%
  \end{tcolorbox}%
}

% --- Butir checklist --------------------------------------
\newenvironment{checklist}{%
  \begin{itemize}[label={\color{ubGreen}\large$\square$},
    leftmargin=1.6em, itemsep=0.35em, topsep=0pt]
}{%
  \end{itemize}
}

% --- Peta kemampuan sesi ----------------------------------
\newcommand{\petasesi}[3]{%
  \par\noindent
  {\small
  \renewcommand{\arraystretch}{1.25}
  \begin{tabularx}{\linewidth}{@{}>{\raggedright\arraybackslash}p{0.30\linewidth}
      >{\raggedright\arraybackslash}p{0.33\linewidth}
      >{\raggedright\arraybackslash}X@{}}
    \toprule
    {\footnotesize\sffamily\color{ubMuted}SEBELUM SESI INI} &
    {\footnotesize\sffamily\color{ubAccent}DITAMBAHKAN} &
    {\footnotesize\sffamily\color{ubMuted}MENYUSUL} \\
    \midrule
    #1 & #2 & #3 \\
    \bottomrule
  \end{tabularx}}
  \par\vspace{0.6em}
}

% --- Daftar periksa pada tabel ----------------------------
\newcommand{\centang}{\color{ubGreen}$\square$}
